% Chương 3. Mỗi mục ghi người phụ trách và file nháp để chép nội dung sang. Quyết định đã chốt: docs/03_analysis_protocol.md.
\chapter{Methods}
\label{ch:methods}

\section{Pipeline overview}
\fillpara{[R1] Một đoạn mô tả toàn bộ pipeline theo hình \ref{fig:pipeline}: dữ liệu thô, tiền xử lý, epoch, hai nhánh classical và DL, đánh giá, giải thích.}

\begin{figure}[H]
\centering
\reportfig{F1_pipeline.png}
\caption{Analysis pipeline. \fillin{chú thích}}
\label{fig:pipeline}
\end{figure}

% ------------------------------------------------------------------
\section{Preprocessing}
\label{sec:preprocessing}
\fillpara{[R2] Nguồn chính: \src{reports/report_template.md} mục 3.1 (đã viết), \src{reports/r2_methods_v2.md}. Chia theo các mục nhỏ dưới đây.}

\subsection{Channels and electrode positions}
\fillpara{[R2] 127 kênh EEG, không thêm FCz; vị trí điện cực CapTrak theo từng subject, trường hợp dùng file .bvef; bỏ kênh không phải EEG.}

\subsection{Bad channels and recording joins}
\fillpara{[R2] Tiêu chí phát hiện kênh hỏng (robust z-score, ngưỡng kênh phẳng), nội suy spherical spline; đánh dấu BAD quanh chỗ nối recording.}

\subsection{Filtering and re-referencing}
\fillpara{[R2] Lọc FIR pha không 0,1 đến 40 Hz, lý do không dùng notch, tham chiếu trung bình.}

\subsection{Independent component analysis}
\fillpara{[R2] Extended Infomax, dữ liệu lọc 1 Hz, giữ 99\% phương sai, seed; ICLabel, ngưỡng loại IC mắt và cơ; số IC loại trung bình mỗi người; hạn chế (ICA fit trên toàn bộ recording).}

\subsection{Epoching and rejection}
\fillpara{[R2] Resample 250 Hz, epoch -200 đến 800 ms, baseline -200 đến 0 ms, ngưỡng peak-to-peak, loại epoch chồng BAD, tiêu chí loại subject theo tỉ lệ target giữ lại.}

\subsection{Quality control}
\fillpara{[R2] Các kiểm tra tự động khi đọc lại file FIF, kiểm tra trực quan cho 35 subject, đối chiếu với events gốc; những gì chưa kiểm tra được (sub-10, sub-13 candidate).}

% ------------------------------------------------------------------
\section{ERP analysis}
\label{sec:erp-methods}
\fillpara{[R3] Grand-average ERP tại Pz, CPz, Cz (mean và SEM qua subject); difference wave target trừ standard; topomap 250 đến 600 ms và các lát cắt 300, 400, 500, 600 ms; ghép kênh theo tên điện cực. Quy tắc đo đỉnh P300 theo subject (cửa sổ, kênh, cách xử lý đỉnh ở biên). Nguồn: \src{reports/r3_review_and_handoff.md}.}

% ------------------------------------------------------------------
\section{Classification tasks and time windows}
\label{sec:tasks}
\fillpara{[R1] Bài toán chính và hai bài toán đối chứng, vì sao cần chúng; hai cửa sổ thời gian và lý do chọn 450 ms (RT trung bình trừ 1,5 SD). Nguồn: protocol mục 2 và 3.}

\begin{table}[H]
\centering
\caption{Classification tasks and time windows. \fillin{chú thích}}
\label{tab:tasks}
\begin{tabularx}{\linewidth}{llX}
\toprule
\textbf{Name} & \textbf{Definition} & \textbf{Purpose} \\
\midrule
\code{tgt_std} & \fillin{} & \fillin{} \\
\code{dis_std} & \fillin{} & \fillin{} \\
\code{tgt_dis} & \fillin{} & \fillin{} \\
\midrule
\code{full} & \fillin{} & \fillin{} \\
\code{early} & \fillin{} & \fillin{} \\
\bottomrule
\end{tabularx}
\end{table}

% ------------------------------------------------------------------
\section{Classical models}
\label{sec:classical}
\fillpara{[R4] Nguồn: \src{reports/methods_ml_rq2_draft.md} mục 3.3.}

\subsection{Time-window shrinkage LDA}
\fillpara{[R4] \code{lda_tw}: trung bình biên độ theo bin 50 ms trên mọi kênh, số đặc trưng cho mỗi cửa sổ, chuẩn hoá, shrinkage LDA với prior bằng nhau.}

\subsection{xDAWN covariances and Riemannian tangent space}
\fillpara{[R4] \code{xdawn_rg}: số bộ lọc xDAWN, ma trận hiệp phương sai, tangent space, logistic regression có class weight; vấn đề hạng thấp sau ICA và cách xử lý.}

\begin{table}[H]
\centering
\caption{Classical model settings. \fillin{chú thích}}
\label{tab:classical-params}
\begin{tabularx}{\linewidth}{lXX}
\toprule
\textbf{Model} & \textbf{Features} & \textbf{Classifier and hyperparameters} \\
\midrule
\code{lda_tw} & \fillin{} & \fillin{} \\
\code{xdawn_rg} & \fillin{} & \fillin{} \\
\bottomrule
\end{tabularx}
\end{table}

% ------------------------------------------------------------------
\section{Deep learning models}
\label{sec:dl}
\fillpara{[R5] Nguồn: \src{reports/methods_dl_draft.md}.}

\subsection{Architectures}
\fillpara{[R5] EEGNet-8,2 và ShallowFBCSPNet: các lớp chính, số tham số; viết bằng PyTorch, cùng kiến trúc với braindecode.}

\subsection{Input and normalisation}
\fillpara{[R5] Resample cả epoch về 128 Hz rồi mới cắt cửa sổ; z-score theo kênh với thống kê chỉ từ phần train.}

\subsection{Training}
\fillpara{[R5] Optimizer, learning rate, batch size, số epoch tối đa, early stopping trên validation tách từ train, class weight, seed, phần cứng và thời gian chạy.}

\begin{table}[H]
\centering
\caption{Deep learning training settings. \fillin{chú thích}}
\label{tab:dl-params}
\begin{tabular}{lcc}
\toprule
\textbf{Setting} & \textbf{EEGNet} & \textbf{ShallowFBCSPNet} \\
\midrule
Number of parameters & \fillin{} & \fillin{} \\
Optimizer, learning rate & \fillin{} & \fillin{} \\
Batch size & \fillin{} & \fillin{} \\
Maximum epochs, patience & \fillin{} & \fillin{} \\
Validation split & \fillin{} & \fillin{} \\
\bottomrule
\end{tabular}
\end{table}

% ------------------------------------------------------------------
\section{Evaluation}
\label{sec:evaluation}

\subsection{Within-subject and cross-subject validation}
\fillpara{[R4, R5] Within: 5 đoạn liên tiếp đều nhau, không xáo trộn (xem changelog protocol 2026-11-19). Cross: GroupKFold 5 fold theo subject, file fold cố định dùng chung cho mọi mô hình.}

\subsection{Leakage safeguards}
\fillpara{[R4] Mọi bước học từ dữ liệu chỉ fit trên train; tuning trên validation tách từ train; cân bằng lớp bằng class weight; seed cố định.}

\subsection{Metrics}
\fillpara{[R4] Balanced accuracy và ROC-AUC là metric chính, F1 lớp target là phụ; mức ngẫu nhiên 0,5; vì sao không dùng accuracy thô.}

\subsection{Statistical comparison}
\fillpara{[R4] Đơn vị thống kê là subject; chọn classical tốt nhất và DL tốt nhất theo trung vị balanced accuracy cửa sổ full; Wilcoxon signed-rank hai phía, hiệu chỉnh Holm, rank-biserial; công cụ (pingouin). Nguồn: protocol mục 8.}

% ------------------------------------------------------------------
\section{Model explanation}
\label{sec:explain-methods}
\fillpara{[R1] Nguồn: \src{reports/explain_dl_draft.md} mục 3.6. Chia theo các mục nhỏ dưới đây.}

\subsection{EEGNet spatial patterns}
\fillpara{[R1] Chuyển trọng số depthwise thành activation pattern, chuẩn hoá về µV, quy ước dấu, tương quan với bản đồ P300; độ quan trọng của filter đo bằng mức giảm AUC khi bỏ filter.}

\subsection{Integrated Gradients}
\fillpara{[R1] Target của IG (logit difference), baseline, số bước, kiểm tra completeness, trial được giải thích, cách lấy trung bình theo subject.}

\subsection{LDA activation patterns}
\fillpara{[R1] Fit lại \code{lda_tw} trên cùng 5 fold cross-subject, công thức pattern theo \citet{haufe2014}, kiểm tra kết quả fit lại trùng với R4.}

\subsection{Regional summary and control tasks}
\fillpara{[R1] Hai vùng ROI (P300 và vận động trái), 4 bin thời gian, chỉ số enrichment, kiểm định Wilcoxon mô tả; huấn luyện EEGNet trên hai task đối chứng với cùng code, fold, hyperparameter.}

% ------------------------------------------------------------------
\section{Implementation}
\fillpara{[R1] Ngôn ngữ, thư viện chính và phiên bản (MNE-Python \citep{gramfort2013}, scikit-learn, pyriemann, PyTorch, pingouin), phần cứng, cách chạy lại (lệnh scripts theo thứ tự số).}
